\section{A functional account of consciousness}
\label{sec:account}
\label{sec:approach}
% Frames A4 (definition), A1--A2 (function), A3 (lamp), D1 + D12 + A6c (grammar, precedent, criterion).

\subsection{The definition}
\label{sec:definition}
% A4

Take an agent that acts in an environment, refers to itself, and reasons about itself, either
explicitly or through a proxy such as reasoning about its own utterances. Its consciousness is the part
of its self-reasoning that describes its own reality as it has it: what is there, what it is like, and
where the agent stands among the things described.

Three properties of the definition matter later. It defines a function and not a substance, so the
question ``how much'' is well formed. It says nothing about substrates, on purpose. And it is not a
criterion one can apply by inspection: whether a given system's self-reasoning has this part, and how
much of it, is what the rest of the paper is about.

\paragraph{Batch and incremental forms.}
\label{sec:incremental}
The definitions of this section, and the derivations of Section~\ref{sec:model}, are stated in batch
form: a function takes an input and returns a result, and the Observer will be derived as a
conclusion, a discrete output. Human consciousness presents itself as continuous, and the Observer in a
human is stretched over time rather than delivered once. The difference is one of form and not of
content. Any computation stated in batch form has an incremental form, in which the result is
maintained under changes to the input from the previous result and the change, instead of being
recomputed from the whole input, and the transformation from one form to the other is a studied,
systematic one \citep{liu2025}; the RETE algorithm is the incremental form of rule matching
\citep{forgy1982}, and the same transformation applies to any function of this paper. The batch form is
used below because it is the one in which the derivations are shortest. The incremental form is the
one in which the functions run, in a brain by frames and in a Transformer by layers and tokens, and
Section~\ref{sec:constructions} states what the transformation has to deliver: the conclusion
re-derived every frame at low cost from the previous conclusion and what the frame adds, so that
continuity is the run of re-derivations and nothing else.

\subsection{A function is what recurs, and what recurs is compressible}
\label{sec:function}
% A1, A2

A function, in this paper, is anything that recurs across occasions: a piece of the agent's processing
that is the same on many occasions of its running. The definition is deliberately thin: it mentions no
purpose, no design and no mind, so that nothing about consciousness is smuggled in at the start, and the
recurrence is across time, not the internal regularity of a single structure. With
Section~\ref{sec:definition} it composes: the functions of consciousness are what recurs in the part of
an agent's self-reasoning that describes its own reality.

Recurrence has an information-theoretic consequence, and the thin definition inherits it. If a pattern
is produced often by a computable process, it has high algorithmic probability, in Solomonoff's sense
of the probability that a universal machine fed random bits outputs it
\citep{solomonoff1964a,solomonoff1964b}. By the coding theorem, $K(x) = -\log m(x) + O(1)$: an object
with high algorithmic probability has a short description, up to an additive constant
\citep{levin1974,livitanyi2019}. So a function in our sense is compressible, and we did not have to
assume compressibility; it came from recurrence. The constant depends on the reference machine, which
is to say on who is doing the compressing, and the gap between the constants of two substrates is what
makes the comparison of a model with a human nontrivial; Section~\ref{sec:profile} turns that gap into
the object of measurement.

Read in the other direction, the same theorem says that a search process finds simple things first,
whether the search is evolution or gradient descent. This has been made quantitative. Input--output
maps of a wide class are strongly biased toward simple outputs, with the bias following an upper bound
of the form $P(x) \le 2^{-a \tilde{K}(x) + b}$ \citep{dingle2018}; the parameter--function map of deep
networks is biased toward simple functions, which is offered as an explanation of why they generalize
at all \citep{valleperez2019}; and stochastic gradient descent lands on functions with probabilities
close to those of a Bayesian sampler over the same prior \citep{mingard2021}. The argument uses both
readings: backward, from recurrence to short description, to define functions; forward, from search to
simplicity, to expect a trained model to have found the ones that recur in its data.

\subsection{Computational functionalism, and the lamp}
\label{sec:lamp}
% A3

Something is lit in a human's report of themselves. It is not lit in the same way at all times, and its
absence in dreamless sleep or under anaesthesia is part of the data. We take the report as data rather
than as error, which is what distinguishes the position here from an eliminativism that treats
self-report as noise.

Consider a physicalist who is not a panpsychist. We agree with them that the lit thing is physical and
requires no second substance, and that there are no abstract machines: everything that runs is a
physical object, and a computation without a physical realization is a description of a computation.
Then a question arises for them. What prevents a physical machine on a non-protein substrate from
having the lit thing? We do not claim the answer is ``nothing''. We note that whoever answers ``the
protein'' owes a proof, and that a proof of that shape would be vitalism: an appeal to a property of a
material that does no work in any other explanation of what the material does. Pending such a proof we
adopt computational functionalism, with as little physics as we can manage
(Section~\ref{sec:physics}).

One objection deserves its own reply, because it is the strongest of the substantialist objections. A
simulated superconductor carries no current: run the simulation and no magnet levitates, because the
current in the simulation is a number and a magnet needs charge. So, the objection goes, a simulated
brain feels nothing. The reply is that the objection needs a consumer of the current outside the
simulation. The simulated superconductor fails because we wanted the current in our substrate, for a
magnet in our substrate. Consciousness has no such external consumer: whatever uses it is the agent,
and the agent runs in the same substrate as the process. Someone who wants the objection to work must
name a function of consciousness whose consumer lies outside the substrate where the agent runs. None
has been named, and the natural candidates, control of the agent's own behaviour, coordination with
other agents, allocation of its own effort, all have consumers inside.

\subsection{Method: change the grammar, and keep a criterion for content}
\label{sec:grammar}
\label{sec:precedent}
\label{sec:criterion}
% D1, D12, A6c

``Consciousness'' is a noun, and a noun invites the questions one asks about things: where is it, what
are its properties, when does it begin, does this system have one. Every sustained attempt to answer
those questions ends in one of a few circles. Who observes the observer. How could red look like
anything other than the way it looks. Either everything has it, down to the electron, or the line
between what has it and what does not falls somewhere that cannot be justified. The circles recur
across two and a half thousand years of work by careful people, which is evidence that the fault lies
in the object, or in the grammar that made an object of it, and not in the people. An account that
answers ``what is consciousness'' with a definition of a thing inherits the circles whatever the thing
is made of.

A science met circles of this kind until it changed the grammar. Gravity was a force, a thing acting
between things, and the attempts to carry it into a relativistic frame as one more field kept failing
in the same places; Nordström's scalar theory was the most careful of them and predicted no light
bending \citep{norton1992}. General relativity dropped the thing. Gravity became the shape of
spacetime, present only in how free bodies move through it, and the move came with a criterion: what
stays the same under every change of coordinates, the curvature, is the physical content, and what can
be transformed away by choosing a frame is bookkeeping. A uniform gravitational field can be removed by
falling; tidal curvature cannot be removed by anything. The geometric idea had been voiced before there
was a theory to carry it \citep{riemann1873,clifford1876}.

We take the same route. Consciousness becomes a shape of reasoning, visible only in how reasoning
moves, and the criterion transfers as follows. An agent can change its self-model in the way an
observer in relativity changes coordinates. Some displacements of its reasoning disappear when the
self-model is refined: they were artifacts of the coarser model, and the agent that noticed them was
reading its own bookkeeping. Others survive every self-model the agent can afford to hold. Those are
the analogue of curvature, and only they count as content; we call them the \emph{computational
curvature} of the agent's reasoning.\footnote{In the first version of this work and in the material
that preceded it, the same displacements were called higher-order computational phenomena (HOCP). The
old name is kept only as an alias for readers of that material.} The criterion turns ``systematic'' from a
description into a test; it gives a procedure for adding a function to the list of
Section~\ref{sec:reference}, exhibit the displacement, then show that refining the self-model does not
remove it; and it rules things out, which a criterion must do to be worth having. An agent's belief
that its decisions are made in the order it remembers them is removed by a better self-model, and it is
therefore coordinates, however vivid it is.

The first version of this work \citep{smirnov2026v1} was written in the vocabulary of qualia, inside the grammar, and met the
wall repeatedly and in the same places. The philosophy of mind stays in the account in two roles: as
the record of where the wall is, and, once the grammar is changed, as a framework that works again for
the named part of the functions (Section~\ref{sec:h1h2}).
